\documentclass[11pt]{article}

%------------------------- fonts & encoding -----------------------------------
\usepackage[T1]{fontenc}
\usepackage[utf8]{inputenc}
\usepackage{amsmath,amssymb,amsthm}
\usepackage{mathtools}
\usepackage[sc]{mathpazo}
\linespread{1.06}
\usepackage{microtype}

%------------------------- layout ---------------------------------------------
\usepackage[a4paper,margin=1.05in,marginparwidth=1.0in]{geometry}
\usepackage{booktabs}
\usepackage{enumitem}
\setlist{itemsep=2pt,topsep=3pt}
\usepackage{xcolor}
\usepackage{graphicx}
\usepackage{tikz}
\usepackage{pgfplots}
\pgfplotsset{compat=1.17}
\usepackage[most,breakable]{tcolorbox}

%------------------------- palette --------------------------------------------
\definecolor{crimson}{HTML}{A51C30}
\definecolor{ink}{HTML}{1A1A1A}
\definecolor{steel}{HTML}{2C5F8A}
\definecolor{amber}{HTML}{9A6A00}
\definecolor{teal}{HTML}{1F6F6B}
\definecolor{paperlt}{HTML}{F7F4EF}

%------------------------- section styling ------------------------------------
\usepackage{titlesec}
\titleformat{\section}{\Large\bfseries\color{steel}}{\thesection}{0.7em}{}
\titleformat{\subsection}{\large\bfseries\color{ink}}{\thesubsection}{0.6em}{}
\titlespacing*{\section}{0pt}{1.6em}{0.6em}

%------------------------- hyperref -------------------------------------------
\usepackage[colorlinks=true,linkcolor=steel,citecolor=steel,urlcolor=crimson]{hyperref}

%------------------------- theorem environments -------------------------------
\theoremstyle{definition}
\newtheorem{definition}{Definition}[section]
\newtheorem{example}{Example}[section]
\newtheorem{exercise}{Exercise}[section]
\theoremstyle{plain}
\newtheorem{proposition}{Proposition}[section]
\theoremstyle{remark}
\newtheorem*{remark}{Remark}

%------------------------- callout boxes --------------------------------------
\newtcolorbox{intuition}{
  enhanced, breakable, colback=crimson!5, colframe=crimson, boxrule=0.5pt,
  left=8pt,right=8pt,top=6pt,bottom=6pt, arc=2pt,
  colbacktitle=crimson, coltitle=paperlt, fonttitle=\bfseries, title={Intuition}}

\newtcolorbox{pitfall}{
  enhanced, breakable, colback=amber!7, colframe=amber, boxrule=0.5pt,
  left=8pt,right=8pt,top=6pt,bottom=6pt, arc=2pt,
  colbacktitle=amber, coltitle=paperlt, fonttitle=\bfseries, title={Pitfall}}

\newtcolorbox{keyresult}{
  enhanced, breakable, colback=steel!6, colframe=steel, boxrule=0.5pt,
  left=8pt,right=8pt,top=6pt,bottom=6pt, arc=2pt,
  colbacktitle=steel, coltitle=paperlt, fonttitle=\bfseries, title={Takeaway}}

\newtcolorbox{observed}{
  enhanced, breakable, colback=teal!6, colframe=teal, boxrule=0.5pt,
  left=8pt,right=8pt,top=6pt,bottom=6pt, arc=2pt,
  colbacktitle=teal, coltitle=paperlt, fonttitle=\bfseries, title={Measured}}

%------------------------- math operators -------------------------------------
\DeclareMathOperator{\Var}{Var}
\DeclareMathOperator{\Cov}{Cov}
\DeclareMathOperator{\Corr}{corr}
\DeclareMathOperator{\diag}{diag}
\DeclareMathOperator{\std}{std}
\DeclareMathOperator*{\avg}{\mathbb{E}}
\newcommand{\T}{^{\!\top}}
\newcommand{\R}{\mathbb{R}}
\newcommand{\half}{\tfrac12}
\newcommand{\wvec}{\mathbf{w}}
\newcommand{\Xmat}{\mathbf{X}}
\newcommand{\Fmat}{\mathbf{F}}
\newcommand{\Vmat}{\mathbf{V}}
\newcommand{\Umat}{\mathbf{U}}
\newcommand{\Rmat}{\mathbf{R}}
\newcommand{\Lam}{\boldsymbol{\Lambda}}
\newcommand{\Del}{\boldsymbol{\Delta}}
\newcommand{\fvec}{\mathbf{f}}

%------------------------- title ----------------------------------------------
\title{\vspace{-2.2em}
  {\color{steel}\rule{\textwidth}{1.3pt}}\\[0.7em]
  \Huge\bfseries Momentum\\[0.35em]
  \large\mdseries The relative-strength factor --- a lagged, exponentially-weighted average of past excess returns\\[0.4em]
  {\color{steel}\rule{\textwidth}{0.7pt}}}
\author{}
\date{}

%------------------------- watermark footer -----------------------------------
\usepackage{fancyhdr}
\pagestyle{fancy}
\fancyhf{}
\renewcommand{\headrulewidth}{0pt}
\renewcommand{\footrulewidth}{0pt}
\fancyfoot[L]{\footnotesize\color{black!45}MSCI Barra USE4 Imitation}
\fancyfoot[R]{\footnotesize\color{black!45}\thepage}

\begin{document}
\maketitle
\thispagestyle{empty}

\begin{abstract}
Momentum is the relative-strength factor: it measures how strongly a stock has been trending, so that recent winners load high and recent losers load low. The empirical regularity it captures is one of the most durable in finance -- stocks that have outperformed over the past several months to a couple of years tend to keep outperforming, and laggards keep lagging -- and USE4 measures it as an exponentially-weighted average of a stock's daily excess returns over a roughly two-year window. Two construction choices define the factor. The first, and the one that distinguishes momentum from a naive trailing return, is a deliberate \emph{lag}: the most recent month of returns is skipped, because at short horizons returns tend to \emph{reverse} rather than continue, and including them would contaminate the medium-term signal. The second is exponential weighting, which lets the more recent (post-lag) returns count for more than the oldest ones. The raw descriptor is a tiny, daily-return-scale number, roughly symmetric but fat-tailed, and it is the least month-to-month stable of the style factors -- a restlessness that culminates in the factor's signature hazard, the occasional violent ``momentum crash'' when leadership abruptly reverses.
\end{abstract}

\begin{tcolorbox}[colback=paperlt,colframe=ink!30,boxrule=0.4pt,arc=2pt,left=10pt,right=10pt,top=8pt,bottom=8pt]
\textbf{How to read these notes.} We move from \emph{why} Momentum is a factor to \emph{how} USE4 specifies it, then to how we actually form the weighted average, with a deliberate stop at the lag that defines the factor (Section~\ref{sec:lag}) and at the cross-section's small scale and notable instability (Section~\ref{sec:reading}). Read those two slowly; they carry the weight of the set. The numbers labelled as measured come from a single run of our personal build.

\smallskip
Four coloured boxes recur: \textcolor{crimson}{\textbf{Intuition}} gives the mental picture, \textcolor{amber}{\textbf{Pitfall}} flags what breaks without the fix, \textcolor{steel}{\textbf{Takeaway}} states a load-bearing result, and \textcolor{teal}{\textbf{Measured}} reports real numbers from the build run.

\smallskip
\textbf{Reference material.}
\begin{itemize}[leftmargin=1.4em,topsep=2pt]
  \item External: Menchero, Orr \& Wang (2011), USE4 Empirical Notes (Appendix A) and Methodology Notes (Sec. 2.2--2.3).
\end{itemize}
\end{tcolorbox}

\tableofcontents
\bigskip

\section{What Momentum is, and why it is a factor}\label{sec:what}

\paragraph{Learning objectives.}
\begin{itemize}[leftmargin=1.4em,topsep=2pt]
  \item State Momentum as recent relative strength -- the tendency of trends to continue.
  \item Explain why trend-following exposure is a shared, co-moving risk.
  \item Say in one sentence what the Momentum exposure is, before any standardisation.
\end{itemize}

Momentum is the observation that trends persist. A stock that has beaten the market over the past year tends, on average, to keep beating it for a while; one that has lagged tends to keep lagging. This medium-horizon continuation -- distinct from any fundamental valuation -- is among the most robust empirical patterns in asset returns, documented across decades, countries, and asset classes. Momentum is pure price behaviour: it asks nothing about earnings or book value, only how the stock has been moving.

That makes it a genuine risk dimension. Recent winners tend to be similar stocks -- often crowded into the same themes -- and they rise and, crucially, \emph{fall} together. When a long-running trend snaps, the whole high-momentum cohort reverses at once, a coordinated move a risk model must be able to see. Momentum is therefore not just a return-predicting signal but a source of shared risk, and USE4 carries it as a style factor. The exposure is a standardised version of a lagged, exponentially-weighted average of the stock's recent excess returns.

\begin{intuition}
The whole factor answers one question: \emph{how strong has this stock's recent run been, relative to cash?} A high score is a stock on a sustained tear; a low score is a persistent laggard. Hold that picture -- a weighted average of how it has been doing -- and the lag and the decay below are just refinements that make the average measure the \emph{right} stretch of the past.
\end{intuition}

\begin{keyresult}
The Momentum exposure is a standardised, lagged, exponentially-weighted average of a stock's daily excess returns. Everything that follows is about which stretch of the past to average, and how to weight it.
\end{keyresult}

\section*{Notation}
\addcontentsline{toc}{section}{Notation}

\begin{center}
\begin{tabular}{@{}ll@{}}
\toprule
Symbol & Meaning \\
\midrule
$r_{i,\tau}$ & excess daily log return of stock $i$ on day $\tau$ (return over the risk-free rate) \\
$L,\ T$ & lag and window length ($L\approx 21$, $T\approx 504$ trading days) \\
$\omega_\tau$ & exponential weight, $\exp\!\big(-\ln 2\,(t-\tau)/h\big)$, half-life $h\approx 126$ days \\
$W$ & full-window weight sum (the normaliser) \\
$m_i$ & raw momentum, $\big(\textstyle\sum_\tau \omega_\tau\, r_{i,\tau}\big)/W$ over the lagged window \\
$x_i$ & standardised MOM exposure for stock $i$ \\
$w_i$ & cap weight of stock $i$ in the ESTU \\
$\avg_{\text{cw}}, \std_{\text{ew}}$ & cap-weighted mean, equal-weighted standard deviation over the ESTU \\
$Q_1,\dots,Q_5$ & quintiles of raw $m_i$, from losers ($Q_1$) to winners ($Q_5$) \\
ESTU & estimation universe (MSCI USA IMI proxy) \\
\bottomrule
\end{tabular}
\end{center}

\section{How USE4 defines it}\label{sec:def}

\paragraph{Learning objectives.}
\begin{itemize}[leftmargin=1.4em,topsep=2pt]
  \item State the window, the lag, the half-life, and what is averaged.
  \item Note the excess-log-return inputs and the standardisation target.
\end{itemize}

USE4 defines Momentum (its relative-strength descriptor) as the exponentially-weighted mean of a stock's daily excess log returns over a trailing two-year (504 trading-day) window, lagged by one month (21 trading days), with a six-month (126 trading-day) half-life (Menchero, Orr \& Wang 2011, Empirical Notes, App. A). Excess means in excess of the risk-free rate; returns are logarithmic, so they add cleanly over time; and prices are split- and dividend-adjusted, so corporate actions do not masquerade as performance.

The lag is not incidental -- it is central, and Section~\ref{sec:lag} is devoted to it. The resulting raw momentum is then trimmed at three standard deviations (Methodology Notes, Sec. 2.2) and standardised over the estimation universe to cap-weighted mean $0$ and equal-weighted standard deviation $1$ (Sec. 2.3). Like Beta, Momentum is built from a stock's own return history rather than read off a filing, and every input is point-in-time.

\section{How we compute it: a lagged, weighted average}\label{sec:compute}

\paragraph{Learning objectives.}
\begin{itemize}[leftmargin=1.4em,topsep=2pt]
  \item Write the weighted-average formula and identify the window and the lag.
  \item Explain the normalisation and why it attenuates short histories.
\end{itemize}

For each stock at each signal date we form a weighted average of its daily excess log returns over the window, skipping the most recent $L$ trading days:
\[
m_i \;=\; \frac{1}{W}\sum_{\tau} \omega_\tau\, r_{i,\tau},
\qquad
\omega_\tau = \exp\!\big(-\ln 2\,(t-\tau)/h\big),
\qquad
\tau \in [\,t-L-T,\; t-L-1\,],
\]
where the sum runs over the lagged window -- from about one month ago back about two years -- and $W$ is the sum of the weights over a full window. The exponential weight gives a return at the recent (post-lag) edge of the window the most influence, with influence halving every half-life of about six months as we look further back, fading smoothly to almost nothing at the two-year end (Figure~\ref{fig:mom_dist}, left).

Dividing by the full-window weight sum $W$ does two things. It turns the quantity into a weighted \emph{average} daily return -- a number on the scale of a single day's return -- rather than a cumulative sum. And it \emph{attenuates} a stock with an incomplete history: a young listing that fills only part of the window is divided by the full $W$ anyway, so its momentum is shrunk toward zero in proportion to the coverage it is missing, rather than being rescaled up to look like a fully-observed stock. That keeps thin-history names from dominating the tails. A stock needs roughly a year of returns within the window to be scored at all; shorter histories are left out.

\begin{pitfall}
Use the right returns. They must be \emph{log} returns (so the window's days add rather than compound awkwardly), in \emph{excess} of the risk-free rate (so a high-rate era does not inflate every stock's momentum), and computed from \emph{adjusted} prices -- a raw price series would read a $2$-for-$1$ split as a $-50\%$ day and brand a healthy stock a loser. And momentum must be normalised by the full-window weight sum, not the observed one: rescaling by the observed weight would extrapolate a six-month run to the two-year horizon and let recent listings explode into the tails.
\end{pitfall}

\section{Why skip the most recent month}\label{sec:lag}

\paragraph{Learning objectives.}
\begin{itemize}[leftmargin=1.4em,topsep=2pt]
  \item Explain the difference between short-horizon reversal and medium-horizon momentum.
  \item Justify the one-month lag as protecting the signal from reversal.
\end{itemize}

This is the choice that makes momentum momentum. Returns behave differently at different horizons. Over the medium term -- several months to a couple of years -- they \emph{continue}: winners keep winning. But over the very short term -- the most recent few weeks -- they tend to \emph{reverse}: last month's biggest winners disproportionately give some back, and last month's worst losers bounce. This short-term reversal has its own well-documented causes, from bid-ask bounce and price-pressure from large trades to simple overreaction that quickly corrects.

If the momentum window ran right up to the signal date, it would mix these two opposing effects: the continuation it wants to capture, and the reversal lurking in the most recent weeks. The reversal would partly cancel the momentum, blunting and adding noise to the signal. The fix is a deliberate lag -- skip the most recent month and measure momentum from there back -- so the descriptor is built only from the horizon where returns continue, cleanly separated from the horizon where they reverse. The gap at the front of the weighting kernel in Figure~\ref{fig:mom_dist} is exactly this skipped month.

\begin{intuition}
Think of two clocks running in opposite directions. The fast clock -- this past month -- tends to \emph{undo} itself: hot stocks cool, cold stocks rebound. The slow clock -- the prior one-to-two years -- tends to \emph{keep going}. Momentum wants only the slow clock, so it covers the fast one's face. Reading both at once would have them fight, and the signal would blur.
\end{intuition}

\section{Reading the cross-section: small, restless, and crash-prone}\label{sec:reading}

\paragraph{Learning objectives.}
\begin{itemize}[leftmargin=1.4em,topsep=2pt]
  \item Interpret the tiny scale of the raw descriptor and read it correctly.
  \item Explain why momentum is the least stable factor and what a momentum crash is.
\end{itemize}

Two features of the cross-section are worth meeting directly. First, the \emph{scale}: because the raw descriptor is a weighted \emph{average} of daily returns, it is a tiny number -- a median around four ten-thousandths, i.e. a few hundredths of a percent per day (Figure~\ref{fig:mom_dist}, right). That smallness is the unit, not weakness; a stock with raw momentum of $0.002$ has been compounding a healthy edge day after day. As always, the standardised score, not the raw number, is the cross-sectional quantity you compare. The distribution is roughly symmetric, slightly left-leaning, and distinctly fat-tailed -- a population of extreme winners and extreme losers beyond what a bell curve would predict.

Second, the \emph{instability}. Momentum is the least month-to-month stable of the style factors. This is structural, not a defect: the window rolls forward each month, the lag edge advances, and old strong (or weak) months drop off the back, so the membership of ``recent winners'' churns far more than a slow-moving balance-sheet quantity ever would. And that restlessness has a violent tail: momentum is subject to occasional \emph{crashes}, episodes -- the recoveries of 2001, 2009, and 2020 are the textbook cases -- where beaten-down losers rebound hardest and the long-winners/short-losers posture reverses sharply. These crashes are why momentum's relationship to the market looks weakly negative over a full sample yet turns neutral once the crisis rebounds are excluded.

\begin{observed}
On the build run the raw momentum descriptor has an ESTU median of $\approx 0.0004$ -- a weighted-average daily excess return of a few hundredths of a percent -- and is roughly symmetric (slight left skew) with fat tails (excess kurtosis $\approx 3$). The three-sigma trim clips $\approx 1.7\%$ of ESTU rows. After standardisation the cap-weighted mean is $0$ to machine precision and the equal-weighted standard deviation is $1.0$; $\approx 98\%$ of ESTU rows fall inside $[\pm 3]$, but only $\approx 93\%$ of all scored stocks do, since the standardisation is calibrated on the investable universe. Quintiles of raw momentum are strictly monotone, from a clear loser quintile to a clear winner quintile. Month-over-month rank stability is the lowest of the style factors, Spearman $\rho \approx 0.87$, a direct consequence of the rolling window. As an economic check, sorting the universe by momentum gives strictly monotone forward returns -- a top-minus-bottom quintile spread on the order of $+1\%$ per month (roughly $+14\%$ annualised) over the sample -- the momentum premium that makes the factor real, even though it is punctuated by the crashes above. The deliverable is $\approx 1.8$M scored rows over $\approx 329$ dates and $\approx 16{,}050$ tickers.
\end{observed}

\begin{figure}[htbp]
\centering
\includegraphics[width=\linewidth]{figures/mom_dist.png}
\caption{Momentum, from the build. \emph{Left:} the weighting kernel -- the most recent month is skipped entirely (the shaded lag, where short-term returns reverse), after which each day's weight decays exponentially with a six-month half-life across the roughly two-year window. \emph{Right:} the raw descriptor, pooled over in-ESTU stock-dates, is a tiny weighted-average daily excess return (median $\approx 0.0004$), roughly symmetric (slightly left-leaning) and fat-tailed; the standardised score, not this raw number, is the cross-sectional exposure.}
\label{fig:mom_dist}
\end{figure}

\section{Trim and standardisation}\label{sec:std}

\paragraph{Learning objectives.}
\begin{itemize}[leftmargin=1.4em,topsep=2pt]
  \item State the standardisation target and why the trim handles the fat tails.
\end{itemize}

The raw momentum $m_i$ is trimmed at three standard deviations and standardised over the ESTU to cap-weighted mean $0$ and equal-weighted standard deviation $1$:
\[
x_i \;=\; \frac{m_i - \avg_{\text{cw}}(m)}{\std_{\text{ew}}(m)},
\qquad
\avg_{\text{cw}}(m) = \frac{\sum_i w_i\, m_i}{\sum_i w_i}.
\]
Because the raw distribution is fat-tailed but roughly symmetric, the trim clips a modest fraction from \emph{both} ends -- the extreme winners and the extreme losers alike -- before the standardisation fixes the centre and spread. The cap-weighted centring measures a stock's momentum against the dollar-weighted market; the equal-weighted scaling sets the unit.

\begin{exercise}
Explain, in two sentences, why including the most recent month in the momentum window would tend to \emph{weaken} the factor rather than strengthen it. (Hint: what does the most recent month's return tend to do next?)
\end{exercise}

\begin{exercise}
A stock's raw momentum is $0.0015$ and another's is $-0.0010$. Both are tiny numbers; explain what they mean in words (per-day terms), and why you would still expect a meaningful gap between the two stocks' recent performance over the whole window.
\end{exercise}

\begin{exercise}
Momentum's month-over-month rank stability ($\rho \approx 0.87$) is the lowest of the style factors, while a balance-sheet factor such as leverage is the highest ($\approx 0.99$). Explain the difference in terms of how each factor's underlying inputs change from one month to the next.
\end{exercise}

\section{Summary, and one caveat}\label{sec:summary}

\paragraph{Learning objectives.}
\begin{itemize}[leftmargin=1.4em,topsep=2pt]
  \item Recite the Momentum build from daily returns to standardised exposure.
  \item State the central limitation and what would address it.
\end{itemize}

The build, end to end: take a stock's daily excess log returns over a roughly two-year window, but \emph{skip the most recent month} to avoid short-term reversal; weight the remaining days exponentially with a six-month half-life, so the recent (post-lag) past counts most; average them, normalising by the full-window weight sum so the number is a weighted-average daily return and short histories are attenuated; require roughly a year of data; trim at three standard deviations, clipping both fat tails; and standardise to cap-weighted mean $0$, equal-weighted standard deviation $1$. The output is a small-scale, fat-tailed, restless relative-strength factor.

\begin{keyresult}
Momentum is a lagged, exponentially-weighted average of past excess returns. The two computations that matter most are the one-month lag -- which excludes the reversal-prone recent weeks so the factor measures only the continuation horizon -- and the normalisation, which keeps the descriptor a weighted average and shrinks thin-history names toward zero.
\end{keyresult}

\begin{pitfall}
One honest limitation dominates: momentum carries its risk in its tail. For long stretches it is a steady, well-paid factor, and then -- at sharp market turning points -- it \emph{crashes}, as routed losers rebound and the winners-over-losers posture reverses violently in a matter of weeks. Its own volatility is therefore strongly time-varying, low in calm trends and spiking in reversals, so treating momentum risk as constant understates the danger exactly when it matters. A standard remedy this build does not apply is to scale the signal by its own recent volatility (a volatility-managed momentum), which tames the crashes at some cost to the calm-period premium. A second, smaller point is that momentum is pure price history with no fundamental anchor, so it can quietly load onto crowded trades. The clear levers are a volatility-scaling overlay and crash-aware risk forecasting for the factor.
\end{pitfall}

\begin{remark}
None of this changes the factor's place in the model: Momentum remains the relative-strength dimension, a major axis of style risk alongside size, value, and beta. The caveat is about \emph{time-varying, crash-prone risk}, not about whether the momentum effect is real.
\end{remark}

\begin{center}\small\itshape
References: Menchero, Orr \& Wang (2011), \textit{The Barra US Equity Model (USE4)} --- Empirical Notes (Appendix A) and Methodology Notes (Sec. 2.2--2.3). ESTU treated as an MSCI USA IMI proxy; measured quantities are from a personal build run.
\end{center}

\end{document}
